\section{在一份制度名目之外：辽的四种接口}

辽朝的北面、南面官制常被画成一条干净的分界线：草原一侧由一种办法治理，燕云州县由另一种办法治理。这幅图有助于初学者记住官名，却很容易把真正费力的部分抹掉。官名不是道路，诏书不是翻译，营卫的编组也不会自己把一车粮食送到城门。一个命令若要从上京抵达南京，必须经过能够读写的人、愿意承办的地方首领、可通行的驿路、可征发的马匹与仓储；若其中一环不合，制度的两面便只剩下史书中整齐的栏目。

因而，所谓双重制度首先应当从接口来读。辽廷需要让不同来源的军队在同一场出征中会合，也需要让燕云的赋役、城防与文书进入皇帝可以处理的次序；地方社会则要在这些要求之外保存耕作、婚姻、市场与宗教生活的节奏。两边都不只是被动地接受分类。翻译、书吏、行商、部众首领、城门守卒、寺院住持和能出面担保的家族，正是把分类变成实际行动的人。他们留下的材料极不平等：官修志书偏重中枢安排，宋使记下自己被接待的路线，墓志则保存显贵家庭愿意展示的关系。没有一类材料能单独告诉我们辽的制度在每一村庄如何运行。

\subsection{上京与南京之间：一条路要经过多少种官}

九三八年幽州升为南京，\eventref{LIAO-0938-EVENT-031}{幽州升为南京}，使辽廷取得的并不只是一个便于南望的城名。燕云的城郭、农田、河流渡口和旧有官署必须能同北方的军政网络一起运作，才会成为可用的资源。城墙能在一日内改变守军的旗帜，仓门、账簿和征调关系却有各自的速度。辽廷若要从南部获得粮草、布帛或人力，须依靠熟悉当地度量、地名和赋役习惯的人；若要让北方的军队抵达这一地区，又要处理草料、马匹、关卡、季节和沿路住宿。南京不是“汉地”被装进一个草原国家的盒子，而是几套原有的资源关系被迫重新接线的地方。

《辽史》百官志与营卫志把北面、南面及诸司列为制度条目，给人的印象仿佛每项职责都有固定边界。但志书是后来把长期变化归整后的文本；它保留下的是名称与秩序，不是每一次命令在什么地点被谁执行。宋人使辽的记载则反过来特别注意馆舍、宴飨、礼仪与使团经过的城镇，因为这些正是使者任务所需。两种材料放在一起，可以看到国家为远行和接待预留了程序，却不能从中推断普通居民每天怎样与官署打交道。城址与墓志提供了另一种尺度：前者让城门、道路和手工业区不再只是文字中的背景，后者偶尔记下任官、祖籍、婚姻和葬地。它们仍然偏向保存较好的地点和有能力立碑的家庭。

双重制度并非把人永久锁在两类社会之中。一个在南京任事的官员可能需要理解南部文书，也可能同北方军政发生联系；一名掌管运输的中介不必在任何志书中占有醒目的位置，却可能决定军粮是否赶在河道封冻前送达。相反，一纸来自上京的指令即使有明确的官署名称，也可能因缺少马匹、仓储或地方合作者而停在半路。我们难以逐件重建这种失败，因为留下来的往往是颁令、褒奖或大事，而不是没有按时抵达的一车草料。正是在这种沉默中，不能把北面、南面理解为两架从不打结的机器。

辽的多都城格局使问题更加具体。不同中心并不表示皇帝可以随意把资源从一处挪到另一处；每一次移动都要承担距离的成本。上京附近的草场、东北的山林河谷、燕云的农业区和南京的市场，对同一项征发给出的条件并不相同。朝廷要调兵时，营卫的关系可能更有用；要维持一座南部城市，州县的文书与地方仓储又不可少。国家能力因此不是一个静止的“强”或“弱”，而是能否在不同季节、不同区域，把这些不相同的条件暂时拼合起来。

这种拼合也让地方家庭看见国家。对城外耕作者而言，政权并不总以皇帝的名号出现，而会以修城的劳役、运粮的催办、田地附近的驻军、市场中改变的查验或亲属被征去服役的消息出现。对城市中的匠人和商贩而言，城防、道路安全与驻军消费可能带来机会，也可能意味着被征取物资或被限制出入。材料不足以让我们给这些人安排统一的态度；他们的处境会随地点、职业、亲属资源和政治局势而异。将他们一概写成“受益的汉人”或“被统治的居民”，反而会重复官修分类的盲点。

\subsection{边市的早晨：和平也要人力、文书和风险}

澶渊之后的边界最容易被误读为一条停战线。其实，\eventref{LIAO-1005-EVENT-038}{辽方落实澶渊盟约}所建立的，是一套经常需要操作的关系：使节要被迎送，银绢要被交收，边防仍须驻守，互市的日期、地点与人员要被不断确认。对两朝中枢而言，盟约使战争成本可以转化为可预计的支出和仪式；对边地而言，它把不确定性压缩为另一类工作。货物到达前，要有道路能走，渡口能过，守卒知道何人可通行，翻译能让对方的文书变成可执行的意思。任何一项失灵，都可能把一场本应例行的交往重新变成边警。

边市不应被浪漫化为不受国家干预的自由贸易。开市本身就是政治决定，货物、人员与消息在关口被检查，城镇与驻军也会从交换中取得利益。可是，它也不能只被缩为“岁币的附属”。沿边的车夫、牲畜贩运者、手工业者、旅舍主人、通事以及为使团供应食物的人，要在开市、闭市和道路状况之间安排自己的生计。对一户依赖市场出售布帛、盐货或畜产的人而言，一次禁令或一段路上的不安全，可能比开封或上京的外交辞令更先改变一年的打算。我们不知道他们多数人的姓名，也不能用后世的物价或民族标签替他们说话；但正是这种依赖，使和平从条约变成地方生活的一部分。

宋辽使节往来所留下的行程文字，能让这一过程偶尔变得可见。使者要报告所见的城镇、仪注与接待，因为朝廷需要知道对方怎样安排自己的代表。这类文字特别善于保存正式接触，较少记录未入仪式的交易与劳动。它可以告诉我们一支使团怎样被领过关口，却不能据此断定周围每个村庄都因盟约兴旺；它也会把陌生风俗写得格外显眼，使读者误以为差异就是全部。将使行与城址、钱币、墓志和制度材料并读，能够提出更节制的问题：哪些地点被反复用作交通节点，哪些身份必须在两套礼仪间解释自己，国家如何借着接待与查验获得信息。

银绢的移动尤其显示了财政与空间不可分开。朝廷账簿可以把它记为一次数额，实际交收却要经过仓储、封识、护送和验收。宋方要在自己的财政体系中筹集与转运，辽方也要把接收的物品放进本朝的礼制、赏赐和资源分配。我们不能从盟约文字推出每一匹绢的去处，更不能把数字改写成一方“买和平”、另一方“受施舍”的简单道德故事。能够较为确定的是，固定交往让双方都必须建设并维护一串中间节点；这些节点既降低大规模战争的频率，也把负担分散到仓吏、役夫、护送者和沿路社区。

边市还使边界变得更有弹性，而不是更不真实。战争时期，越境可能意味着掠夺、逃亡或军队穿插；停战后，同一条路可能成为使者、商队或归还人员经过的通道。人们知道哪里需要文书、何处有驻军、哪个渡口在某一季节更容易通过，于是开始按边界来安排时间。对权力较弱的人，这种秩序并不总是保护：一旦市场关闭或军情紧张，他们往往没有足够资源等待下一次开市。但持续交往也使单纯的敌我叙事难以成立。边地社会既要防备对方，又要依赖对方的道路、货物和消息。

\begin{debate}{澶渊的银绢能否说明辽宋之间的上下关系}
宋、辽史书对盟约的措辞各有本朝礼制与政治记忆，后来的“岁币”“纳贡”之类概括又常把这种语言压成一方高、一方低的图式。它们足以说明双方都把交收当作重要政治程序，却不足以替边市、使节、守备和地方征发规定单一含义。把银绢放回仓储、道路与接待中，读者才能看到：关系的稳定不是一个称谓完成的，而是由许多会被中断的手续维持的。
\end{debate}

\subsection{三方谈判：辽不是宋夏关系的背景}

一〇四四年宋夏议和前后，\eventref{LIAO-1044-EVENT-040}{辽的外交姿态影响三方关系}。若只从宋朝朝廷的焦虑来看，辽似乎是一种可以随时施压的北方力量；若只从辽的本纪看，又容易把西夏写成宋辽边缘的一项事务。实际上，三方各自面对不同的道路与资源。兴庆府及河套、河西的水利和城镇，要支撑西夏的军政；宋朝要顾及陕西、河北的防务与财政；辽廷则要评估南部边境、使节礼仪和自己能够调动的军政资源。每一方的谈判都在试探另两方能否承受新的压力，因而没有谁只是他人的布景。

三方关系的可操作性依赖信息，而信息很少完整。使者携带的国书、边将送出的急报、商旅带来的消息和各地观察到的军队移动，抵达中枢的速度与可信度不同。朝廷可以在文书中宣布一种姿态，却不能立即知道对方的地方官、部众首领或沿路守卒是否会用同样方式理解它。辽的外交通常被写成皇帝与使节的往来，但在实际空间中，还包括让来使安全抵达、为其供给、安排翻译、控制其所能观察的事物，以及从其言行中判断对方内部情形。外交不是悬在城市上空的礼仪，而是一种利用道路、人员和时间差的政治技术。

辽与高丽的关系也提醒我们，辽的外部世界并非只在宋辽边界。册封、来使、军事压力与海上、陆上交通在不同阶段交错；高丽朝廷有自己的王统、礼仪与对北方局势的判断，不是辽廷命令的回声。《高丽史》与辽、宋材料可能在同一交涉上强调不同的时点和称谓，这种差异不能靠选一部“更正确”的正史消除。它迫使读者问：谁在什么政治场合记录了这次交往，谁被记录为来朝或受册，谁的港口和道路被默认在叙事之外。礼物和使节的出现能证明一种正式接触，不能自动量出民间海贸的规模或一方对另一方的实际控制。

东北的女真关系更能揭示外交的地方性。辽廷所面对的不是一块可以任意涂色的边地，而是河谷、山林、驻地、贡纳、互市和部众关系组成的许多小尺度政治空间。九百九十一年宋廷拒绝女真请伐辽，\eventref{SYD-991-LIAO-019}{宋廷不许女真请战}，并不意味着东北竞争就此静止；它只是说明宋朝当时不愿把某一种地方压力直接接到自己的战争选择上。女真行动者也不是后来金朝的预告片，他们要处理的是当时可以取得的物资、盟友与消息。把后来的胜利倒灌回早期交涉，会把他们在不同地点作出的犹豫、联合与冲突全部写成命中注定。

从辽廷的角度说，维系这些外部关系并不等于占有它们。朝贡、册封、贡纳与互市都是可以被争论、被暂停、被重新安排的程序；它们往往说明双方承认对方拥有某种需要处理的力量，未必说明一方已经把另一方纳入同一行政体系。尤其在山地、河谷和海路上，地方首领、商人和使节中介能否合作，常比国书里的名词更快决定关系会不会继续。官修史喜欢把这种复杂性收束为“服”“叛”“朝”“拒”，而道路上的人必须面对的是下一批马匹、下一处渡口和下一次可能改期的会面。

\subsection{谁为一座城交税：地方家户不是制度图的留白}

辽的地方家户很难进入连贯的国家叙事，并非因为他们没有被制度触及，而是因为保存下来的材料很少以他们为中心。墓志可以记录一个家族怎样把籍贯、官历、婚姻和葬地编成值得传后的故事；寺院碑刻可以记下施主、功德与重修；城址与墓葬能显示部分地区的居住、生产和丧葬安排。它们让我们看见若干家庭怎样借任官、迁居、捐施或亲属关系获得位置，却不能为辽境每一户制作社会表格。没有墓志的农户、雇工、役夫、妇女和儿童并不因此没有历史，只是更少在能够留存的文字中拥有自己的句子。

也正因为这样，讨论赋役时不能把食货志中的项目直接换成每户的负担。一个制度条目说明朝廷如何命名收入、征发或仓储；它不告诉我们具体地方是否按期缴纳，地方官怎样折算，谁替谁担保，或灾年之后欠额怎样处理。燕云农业区、东北的城镇与山林周边、上京附近的草原地带，面对的资源结构不同，国家命令抵达的方式也不同。统一的国号不取消这些差异。对家庭来说，登记可能带来保护、差役、运输义务或某种可被援引的资格；哪一种更突出，要看他们处在什么地点、能借助哪些亲属和中介，不能预先由族名决定。

寺院是观察这一问题的一处窗口，却不应被夸大为避开国家的世外空间。佛寺的修建、经卷的抄写、造像与法会，都需要木材、纸张、工匠、施主和地方许可。显贵捐施者可以通过碑文把官衔与功德放在同一块石头上，地方家族也可能借宗教活动巩固可见的关系。这样的碑刻让研究者看见一部分资源如何汇集，却不能证明寺院完全由朝廷操控，更不能证明周围居民都以同一种方式信仰。宗教在这里之所以重要，不是因为它提供了一个脱离经济的“精神世界”，而是因为它把家庭财力、书写技能、地方声望和跨地区网络接到一起。

契丹大字、小字题记与汉文碑志的并存，也不宜被压缩成一张稳定的族群地图。文字选择会随官职、仪式、受众、石材、书写传统和保存机会改变。能读写两种或多种文字的人，可能是制度接口中的关键人物；但从一处双语材料不能推出整个地区的语言能力，更不能直接判断某个人的政治归属。对于本书关心的地方家户而言，这一点格外重要：后世常把他们放进“契丹”“汉人”“渤海遗民”的大类，材料中真正可抓住的却往往是一次迁居、一桩婚姻、一笔捐施或一个官名。大类提供分析起点，不能替代这些关系本身。

环境与季节使家庭和制度的相遇更加不均匀。草场、河水、山口与冰冻决定牲畜、车队和军队能否按预期移动；城市附近的粮食、木材与燃料又依赖远近不同的供应地。本纪中的灾异记录通常是朝廷怎样看待异常的证据，不是全境损失的统计表。我们可以据此追问某一时期交通与征发为何紧张，却不能把一条旱涝记载扩写为人人同样经历的饥荒。材料稀疏处需要保留这种不确定性，才不会用一个宏大的国家故事取消家户在不同环境中承受的差别。

\subsection{一纸文书到一座城门：法律怎样才算抵达}

如果把辽的法律与行政只写成百官志中的官名，最容易失去的是办理本身。一个诉讼、一次税物的催纳、一宗军需的交割或一件田宅继承，首先需要有人把口头争执变成可被官署识别的事项。谁会书写，谁能作保，谁知道该到哪一座城、哪一处官署申诉，谁又有余裕等待答复，都会改变同一项规则的实际形状。中央制度给出的是可援引的语言，城门以内外的人则必须把它翻成担保、证明、封识、传递和催办。

辽代留下的完整基层案卷远少于后世某些地区，因而不能虚构一位普通住户在衙门前的每一句话。可是，这种缺乏本身不应使地方法律从叙事中消失。辽代墓志中对官历和家产的安排、寺院题记中的施主与见证者、以及城市遗址所显示的稳定聚落，都提示人们：身份、财产和义务不是只在皇帝改元时才成为问题。它们每天都要被亲属、邻里、地方首领和官署中介重新确认。我们能说的是，这些关系必然需要某种可通行的程序；不能说的是所有地点都有同一套卷宗、同一种审理速度。

北面与南面制度在这一层面不宜被看成两种互不相干的法律。不同区域会有不同的官号、书写语言和处置惯例，但牲畜争执、运输义务、伤害案件、婚姻与继承并不会因为分类不同而不再发生。一个跨越城镇与草场的纠纷，尤其可能需要翻译、地方担保人和了解多种习惯的人。缺少连续司法档案，使我们难以判定谁在这种协商中更常占优；但也正因为如此，不能把国家志书中清楚的分类误作每个人已经接受的边界。

制度还常在危机中显形。军队经过、边市暂闭、河道阻塞或灾害发生时，原先可以依赖的担保、运输和仓储链条会突然变得紧张。地方官若要完成征调，往往必须和有车马、粮食或亲属网络的人谈判；家户若要避免一次负担，也可能求助于邻里、寺院或能向官署说明情况的中介。这里不是把每个人都写成精明算计的行动者，而是承认他们面对的选择不止“服从”与“反抗”。史料最常保存命令和结果，较少保存这些过程中的讨价还价；叙事必须让这种不对称保持可见。

\subsection{寺院的门槛：功德、仓储与家族声望}

辽代佛教遗存之所以重要，不在于它们让我们给社会贴上一张“笃信佛教”的标签，而在于寺院把许多原本分散的资源集中到了可见的地点。修一座殿、立一通经幢、抄写或雕印经卷，需要木料、纸张、颜料、工匠、运输与供养。施主的姓名、官衔或家族关系被刻在石上，是因为他们希望将这些投入转换为可被后人承认的功德和声望。寺院由此既与超度、礼拜和葬仪有关，也与地方经济、文字能力和家族竞争有关。

这并不表示寺院可以被当作辽代的“民间档案馆”。保存下来的碑刻选择了值得纪念的人与事，往往把困难的筹措过程压缩成完成后的功德；一座寺院获得赐额或贵族施舍，也不能证明它在所有时段都由国家供给。相反，正是题记中反复出现的捐施、重修与见证，使我们看到资源需要不断补充。建筑不会因为一次落成就自动维持，经卷、僧众与来往者也需要持续的物资与关系。把这种持续性放在长时段中看，才不会把宗教遗存误作某一位皇帝意志的石化版本。

寺院还给迁徙者和外来者提供了有限但真实的落脚处。商旅、匠人、官员和朝圣者经过城市与道路时，可能在宗教空间中获得接待、信息或声望联系；不过，现有材料很少记录他们全部的行程。我们不该据此把寺院想象成人人平等的公共旅舍。不同身份的人进入寺院的方式不同，能否捐施、书写、留名或获得保护，取决于财富、关系和地方环境。佛教网络跨越政权与文字的能力，说明边界并非把一切关系隔断；它同样说明跨越边界需要成本，并由少数可被材料看见的人承担。

多语题记在这里尤其具有提醒作用。一段契丹文字、一段汉文或两者并置，可以说明同一仪式空间容纳了多种书写实践，却不能自动翻译成“融合”或“对立”。写作者选择何种文字，可能受官职、预期读者、碑刻格式、宗教传统与书写者能力影响。对研究者而言，最有价值的是让这些选择保持具体：不要让一块石头替整个社会回答身份问题，也不要因为释读困难就退回只读官修史的单线叙事。

\subsection{高丽方向的海与东北方向的河：外交通道不是一根线}

辽朝的外交史往往从宋使入辽的陆路开始，仿佛所有外部关系都由燕云的关口决定。高丽方向提醒我们，东北的河流、沿海港口、山地通道和政治中心也构成不同的交通世界。使节和礼物能够抵达，意味着出发地与接收地之间曾有人准备船只、马匹、食物、翻译和安全安排；它不意味着这条路径全年开放，更不意味着同一路线上的民间往来有同样规模。海上与陆上的道路都受季节、风浪、军情与地方控制的影响。

辽朝、高丽和宋朝的材料各自记录使节时，都有自己的中心。辽廷关心的是来使如何被纳入本朝礼仪；高丽朝廷关心的是王统、边境与对不同强邻的选择；宋人的记录则常把东北消息放进本国北方安全的框架。三者并读，不是为了拼出一篇没有矛盾的“真实外交记录”，而是为了看见每一方希望通过称谓、座次、礼物和到达日期证明什么。某方没有详记一次来使，可能是档案未存，也可能是其政治叙事不以此为要，不能轻易改写成没有交往。

这种外交通道同样影响地方生活。港口附近的供给、边城的接待、道路旁的旅舍和渡口上的查验，都会让一些社区比远离交通线的地方更频繁地遇见外来人员与国家要求。机会与风险往往同时出现：来往可能带来货物、工作与消息，也可能带来征发、警戒和更严格的控制。我们不能从一支使团推断整个地区因外贸而繁荣，却可以看到外交的成本并不是只由宫廷承担。它需要许多没有进入国书的人提供食物、劳力与方向。

当女真诸部的关系在东北日益紧张时，这些道路不再只是交换的载体，也是消息和军队移动的条件。河谷、山口与驻地决定谁能更快集结、谁能切断供给、谁能在某处获得当地支持。后来金朝的形成使人容易把所有道路都写成一条通往胜利的直线；但在事发当时，辽廷、宋廷、女真诸部与高丽方面面对的是彼此不完全可见的选择。道路提供可能性，不提供结局。

\subsection{从马匹到布帛：财政不是一张只在朝廷出现的表}

辽朝的财政很难用一个总数概括。牧地所出、城镇税源、战争缴获、与宋的交收、地方仓储和劳役，并不总能换算成同一种单位；即使史书留下数额，也常服务于奖赏、战功、灾异或某项政策的叙述。把这些数字直接相加，仿佛能够算出国家的富足或衰弱，会遮蔽它们首先来自不同机关、不同地区和不同记录目的。更切近的提问是：某种资源怎样进入可用的仓储，谁能决定它的去处，哪一段运输把成本压在何处。

马匹是一个显而易见的例子。对北方军政而言，马并非抽象的“骑兵优势”，而要有牧养、草料、季节性移动、看护和能够在需要时聚集的关系。战时调马要与道路、饮水和人力同时安排；平时维持马群又需要不同于燕云农田的资源条件。城市和农耕地区提供的布帛、粮食、铁器与文书，不能简单替代这些条件，只能与之交换、补充或被征发。双重制度之所以不是文化标签，正因为它面对的是这种不能互换的资源。

布帛、粮食和金属的流动也有同样的限制。南京或其他南部城市的仓储若要供给驻军，须有人征收、登记、保管和转运；物品在账簿上完成一次归类，并不保证它已经到达需要它的人手中。雨雪、道路损坏、守卒更替、地方抵制和临时军情都可能改变交割。我们没有保存完整的辽代地方仓账，无法精确计算每一段损耗；正因为无法计算，才不应把“南部财赋支持北方”写成没有摩擦的结论。可以确定的是，多中心国家若缺少这种物资协调，任何宫廷与军队的安排都会迅速失去支撑。

市场处在国家财政与家庭生计的交界。商人和工匠的货物会被关津、城门、驻军需求与货币使用影响，但他们并不只是被国家汲取的对象。有人为军队或使团提供运输，有人通过边市找到销路，有人则因禁令、征调或道路中断承受损失。不同人能否把风险转给他人，取决于资产、亲属、行旅网络和地方保护，而不是一个“商人阶层”就能说明。现有材料偏向上层与官方，不足以把这些策略统计为全辽社会的规律；它足以提醒读者，财政落在地方时，是许多不相同的交易、差役和等待。

澶渊后经常交收的银绢可以放回这个问题来读。它既是外交关系中可见的物质承诺，也是两套财政与运输体系不得不反复解决的具体工作。接收的一方如何分配，支付的一方如何筹措，沿途人员如何被雇用或征发，都不可能由盟书本身回答。把银绢理解为“和平的价格”固然简洁，却遗漏了价格由谁支付、怎样被转运、在何处造成新的负担。比起给盟约下一个固定的褒贬，追踪这些无法省略的手续更能看出多政权和平的韧性与脆弱。

\subsection{材料的沉默：不能替辽的住户补写一生}

辽史研究的困难不是完全没有材料，而是材料的保存形状与读者的问题并不相同。帝纪善于排列即位、出征、册封和大事，志书善于归类官署、营卫和制度，宋使行录善于描写被安排看见的道路与礼仪；契丹文字题记、汉文墓志、寺院碑刻与考古遗存则留下少量地点、人名和物质环境。它们彼此照亮，也彼此限制。若只读一种，辽会变成一座由皇帝、部族或“民族”驱动的空城；若硬把它们拼成一套没有缝隙的故事，又会把每种材料的立场和沉默消掉。

因此，本章在写到地方家户时宁可保持克制。一个墓志中的迁居不代表所有人都能移动，一通重修碑不代表寺院覆盖了周围生活，一处城址的繁荣也不等于所有居民富足。反过来，普通人没有留下名字，也不能成为他们没有选择、没有关系或没有承受制度的理由。我们只能从国家命令必须经过道路、仓储、担保、翻译与劳力这一点推知，辽的制度从来不是只在宫廷完成的；至于每一个接口由谁承担、怎样协商，仍有大量问题等待更多文书、释读与区域考古回答。

\subsection{从道宗的长时段到辽末：接口何时失去作用}

一〇五五年道宗即位后，辽进入一个在帝纪上看似平静的长时段。长在位并不意味着制度已经不需要维护；恰恰相反，和平、市场、寺院、城市与东北关系能持续存在，是因为许多没有进入本纪的大量事务仍在被处理。史书在这类时期显得较为安静，容易诱使读者直接跳到一一一四年的女真战争，把此前几十年当作辽朝衰亡的前奏。这样写会失去两个层次：一是人们如何在相对可预期的边境与城市秩序中安排生活；二是这种秩序所依赖的接口为什么在后来会变得脆弱。

辽末的危机不能由一个道德标签解释。宫廷继承、后族和重臣的关系，南部城市及东北边地的资源，面对女真行动时能否调动军队、粮草和信息，都有不同的时间表。一一一四年阿骨打起兵后，\eventref{LIAO-1114-EVENT-043}{东北关系发生断裂}，旧有的贡纳、驻防与部众联络不再保证辽廷能以原有方式获得支持。对辽朝而言，难题不是在地图上失去一块颜色，而是某些人不再替它传递命令，某些道路不再为它运输，某些城镇和部众开始按新的保护与风险重新计算。

这也说明辽亡不是所有地方同时发生的一次事件。燕京的失守会先改变南部城防与交通接口；东北在更早或更晚的时点有自己的军事和政治转折；草原与西部行动者又有不同的路线。耶律大石在一一二四年西行，\eventref{LIAO-1124-EVENT-047}{把辽末政治带入中亚道路}，不是把旧王朝连同全部制度搬到远方，而是带着一部分人、军政经验和政治语言，去寻找新的水源、城镇、盟友与补给。对沿途政权来说，这支队伍也不是抽象的“辽遗民”，而是会改变当地联盟与交通风险的新行动者。

留在原辽境的居民同样不应只用“金代开始”一句带过。新政权要接收的不是空城，而是田地、寺院、作坊、道路、家族与旧有文书关系；地方人要应付的也不只是换一个国号，而是新的征发、保护者、市场通道和可被承认的身份。我们缺乏足以逐户追踪这种转换的连续账册，因此不能把一块墓志或一座寺院重修变成全体辽人的命运。较稳妥的结论是：王朝终结改变了制度接口的归属，却没有使生活在这些接口上的人立即停止安排粮食、亲属、信仰和出行。

辽的制度得失也应在这一层面衡量。多中心与多套官署使它能够把草原、东北、燕云城市与对外使节接入同一政治世界；其代价是，每一处接合都依赖特定人员、道路与资源，不能只靠一项制度名称维持。澶渊后的市场与使节降低了战争的频率，却让财政、守备和地方运输成为持续成本；多语书写、寺院与家族网络扩大了可协作的范围，也使国家无法轻易将所有关系收进单一档案。当东北的军事与地方关系倒转时，正是这些接口的失灵，而不只是帝纪中的一场败战，决定了辽朝为何不能按旧方式继续存在。
